%! Characteristics of Radical Functions — Unit 6, Lesson 6.0 — Slides (Beamer)
\documentclass[aspectratio=169, 11pt]{beamer}
\usepackage{algebra2-beamer}   % provides \forestheader{} and \sectionlabel[color]{}
% Standard coordinate grid + axes: {xmin}{xmax}{ymin}{ymax}
\newcommand{\lgrid}[4]{%
  \draw[step=1, linegray!45, thin] (#1,#3) grid (#2,#4);
  \draw[->, thick, charcoal] (#1,0)--(#2,0) node[below right, font=\tiny]{$x$};
  \draw[->, thick, charcoal] (0,#3)--(0,#4) node[left, font=\tiny]{$y$};}
% Cube root: pgfmath has no cbrt, so plot the two branches separately.
\newcommand{\cbrtpos}[2]{\draw[very thick, forest, domain=#1:#2, samples=70, smooth]
  plot (\x,{pow(\x,1/3)});}
\newcommand{\cbrtneg}[2]{\draw[very thick, forest, domain=#1:#2, samples=70, smooth]
  plot (\x,{-pow(-1*(\x),1/3)});}

\begin{document}

% TITLE SLIDE (CourseName is not defined in beamer, so the name is literal)
{
\setbeamercolor{background canvas}{bg=forest}
\begin{frame}[plain]
  \vspace{2.8cm}\hspace{0.55cm}%
  \begin{minipage}{0.9\textwidth}
    {\color{goldacc}\bfseries\small LESSON 6.0}\\[0.5cm]
    {\color{white}\bfseries\LARGE Characteristics of Radical Functions}\\[1.2cm]
    {\color{forestmid}\small Algebra 2: Shepherd~$\cdot$~Unit 6: Radical Functions}
  \end{minipage}
\end{frame}
}

% HOOK
\begin{frame}[t, plain]
  \forestheader{Two radical functions. Why does one of them stop?}
  \sectionlabel{Hook}
  \begin{columns}[T]
    \begin{column}{0.30\textwidth}
      \[ y = \sqrt{x} \]
      \begin{itemize}\setlength{\itemsep}{4pt}
        \item Begins at the origin and \textbf{never} appears to the left of it.
        \item Which inputs are missing --- and \emph{how many}?
      \end{itemize}
    \end{column}
    \begin{column}{0.34\textwidth}
      \centering
      \begin{tikzpicture}[scale=0.30]
        \lgrid{-4}{9}{-3}{4}
        \begin{scope}
          \clip (-4,-3) rectangle (9,4);
          \draw[very thick, forest, domain=0:9, samples=90, smooth] plot (\x,{sqrt(\x)});
        \end{scope}
        \fill[forest] (0,0) circle (6pt);
      \end{tikzpicture}
    \end{column}
    \begin{column}{0.34\textwidth}
      \centering
      \begin{tikzpicture}[scale=0.30]
        \lgrid{-8}{8}{-3}{3}
        \begin{scope}
          \clip (-8,-3) rectangle (8,3);
          \cbrtpos{0}{8}
          \cbrtneg{-8}{0}
        \end{scope}
      \end{tikzpicture}\\[2pt]
      {\footnotesize $y=\sqrt[3]{x}$ --- runs all the way across}
    \end{column}
  \end{columns}
  \vspace{2pt}
  \begin{block}{Today}
    Every family so far, we asked \emph{what the graph does}. Radical functions make us ask something
    new first: \textbf{where is the graph allowed to exist at all?}
  \end{block}
\end{frame}

% INDEX + RADICAND
\begin{frame}[t, plain]
  \forestheader{The index runs the show}
  \sectionlabel{Definition}
  \begin{columns}[T]
    \begin{column}{0.46\textwidth}
      A \textbf{radical function} puts the variable under a radical.
      \[ f(x)=\sqrt[n]{\;\text{expression in } x\;} \]
      \begin{itemize}\setlength{\itemsep}{4pt}
        \item \textbf{Index} $n$ --- the small number in the notch. A plain $\sqrt{\phantom{x}}$ has
              index $2$.
        \item \textbf{Radicand} --- everything underneath.
      \end{itemize}
    \end{column}
    \begin{column}{0.52\textwidth}
      \begin{block}{From the Warm-Up}
        \centering
        \begin{tabular}{ll}
        $\sqrt{36}=6$ & $\sqrt{-36}=$ \textbf{not real} \\[3pt]
        $\sqrt[3]{8}=2$ & $\sqrt[3]{-8}=-2$ \\
        \end{tabular}
      \end{block}
      \vspace{3pt}
      {\footnotesize \textbf{Even} index $\Rightarrow$ radicand \emph{cannot} be negative --- nothing
      real squares to a negative.\\[3pt]
      \textbf{Odd} index $\Rightarrow$ anything goes --- cubing \emph{keeps} the sign.}
    \end{column}
  \end{columns}
\end{frame}

% FINDING THE DOMAIN
\begin{frame}[t, plain]
  \forestheader{Domain is no longer something you read --- it is something you solve for}
  \sectionlabel[goldacc]{radicand $\ge 0$}
  \begin{columns}[T]
    \begin{column}{0.48\textwidth}
      \begin{block}{The one procedure}
        \textbf{Even} index: set the \emph{radicand} $\ge 0$ and solve. \\[3pt]
        \textbf{Odd} index: domain is all reals. Done.
      \end{block}
      \vspace{4pt}
      {\footnotesize A polynomial was defined \emph{everywhere}. A rational function lost a few
      scattered inputs. A square root loses an entire \textbf{half} of the number line at once.}
    \end{column}
    \begin{column}{0.50\textwidth}
      \begin{itemize}\setlength{\itemsep}{5pt}
        \item $\sqrt{x-3}$: \; $x-3\ge0 \Rightarrow [3,\infty)$
        \item $\sqrt{2x+8}$: \; $2x+8\ge0 \Rightarrow [-4,\infty)$
        \item $\sqrt{5-x}$: \; $5-x\ge0 \Rightarrow (-\infty,5]$ \\
              {\footnotesize \textbf{Flip the sign!} This one opens \emph{left}.}
        \item $\sqrt[3]{x-3}$: \; all reals --- no work needed.
      \end{itemize}
    \end{column}
  \end{columns}
\end{frame}

% SQUARE ROOT PARENT
\begin{frame}[t, plain]
  \forestheader{The square root parent: one arm, and an endpoint}
  \sectionlabel[goldacc]{$y=\sqrt{x}$}
  \begin{columns}[T]
    \begin{column}{0.42\textwidth}
      \centering
      \begin{tikzpicture}[scale=0.34]
        \lgrid{-2}{10}{-2}{4}
        \begin{scope}
          \clip (-2,-2) rectangle (10,4);
          \draw[very thick, forest, domain=0:10, samples=100, smooth] plot (\x,{sqrt(\x)});
        \end{scope}
        \fill[forest] (0,0) circle (6pt);
        \fill[charcoal] (1,1) circle (4pt);
        \fill[charcoal] (4,2) circle (4pt);
        \fill[charcoal] (9,3) circle (4pt);
      \end{tikzpicture}
    \end{column}
    \begin{column}{0.56\textwidth}
      \footnotesize
      \begin{itemize}\setlength{\itemsep}{3pt}
        \item Domain $[0,\infty)$; \; Range $[0,\infty)$
        \item \textbf{Endpoint} $(0,0)$ --- a \emph{solid} dot, because $\sqrt{0}=0$ is perfectly real
        \item Increasing on $[0,\infty)$; never decreasing
        \item \textbf{Absolute minimum} $0$ at $x=0$; \textbf{no} absolute maximum
        \item No symmetry
        \item End behavior: $x\to+\infty \Rightarrow y\to+\infty$ (slowly)
      \end{itemize}
      \vspace{2pt}
      \begin{block}{The honest answer on the left}
        As $x\to-\infty$: \textbf{nothing}. The graph \emph{stops}. That is a complete answer, not a
        missing one.
      \end{block}
    \end{column}
  \end{columns}
\end{frame}

% CUBE ROOT PARENT
\begin{frame}[t, plain]
  \forestheader{The cube root parent: two arms, no edges, no extremes}
  \sectionlabel[goldacc]{$y=\sqrt[3]{x}$}
  \begin{columns}[T]
    \begin{column}{0.42\textwidth}
      \centering
      \begin{tikzpicture}[scale=0.34]
        \lgrid{-9}{9}{-3}{3}
        \begin{scope}
          \clip (-9,-3) rectangle (9,3);
          \cbrtpos{0}{9}
          \cbrtneg{-9}{0}
        \end{scope}
        \fill[charcoal] (-8,-2) circle (4pt);
        \fill[charcoal] (-1,-1) circle (4pt);
        \fill[charcoal] (0,0) circle (4pt);
        \fill[charcoal] (1,1) circle (4pt);
        \fill[charcoal] (8,2) circle (4pt);
      \end{tikzpicture}
    \end{column}
    \begin{column}{0.56\textwidth}
      \footnotesize
      \begin{itemize}\setlength{\itemsep}{3pt}
        \item Domain \emph{and} Range: all reals
        \item \textbf{No endpoint} --- nothing stops the graph
        \item Increasing on $(-\infty,\infty)$
        \item \textbf{No extrema at all} --- falls forever left, climbs forever right
        \item \textbf{Origin symmetry} (odd): $\sqrt[3]{-8}=-2=-\sqrt[3]{8}$, so $f(-x)=-f(x)$
        \item End behavior: both arms, $\to\pm\infty$
      \end{itemize}
      \vspace{2pt}
      {\footnotesize Two arms like an odd-degree polynomial --- but \emph{flattening} as they go,
      not steepening.}
    \end{column}
  \end{columns}
\end{frame}

% COMPARE & CONTRAST
\begin{frame}[t, plain]
  \forestheader{Same unit, two very different pictures}
  \sectionlabel{Compare \& contrast}
  \begin{center}
  \renewcommand{\arraystretch}{1.35}
  \begin{tabular}{lll}
  \hline
   & \textbf{Square root} $y=\sqrt{x}$ & \textbf{Cube root} $y=\sqrt[3]{x}$ \\
  \hline
  Index & even & odd \\
  Domain & $[0,\infty)$ & all reals \\
  Range & $[0,\infty)$ & all reals \\
  Endpoint & yes, at $(0,0)$ & none \\
  Symmetry & none & origin (odd) \\
  Extrema & absolute min $0$ & none \\
  Arms & one & two \\
  Inverse of & $y=x^2$ \emph{(restricted)} & $y=x^3$ \\
  \hline
  \end{tabular}
  \end{center}
  \vspace{4pt}
  \begin{center}
    Every difference on this table traces back to one thing: \textbf{the index}.
  \end{center}
\end{frame}

% INVERSE RELATIONSHIP
\begin{frame}[t, plain]
  \forestheader{Where these two parents came from}
  \sectionlabel[goldacc]{Reflect over $y=x$}
  \begin{columns}[T]
    \begin{column}{0.44\textwidth}
      \centering
      \begin{tikzpicture}[scale=0.30]
        \lgrid{-1}{7}{-1}{7}
        \draw[navy, dashed, thick] (-1,-1)--(7,7);
        \begin{scope}
          \clip (-1,-1) rectangle (7,7);
          \draw[very thick, goldacc, domain=0:2.65, samples=80, smooth] plot (\x,{(\x)*(\x)});
          \draw[very thick, forest, domain=0:7, samples=90, smooth] plot (\x,{sqrt(\x)});
        \end{scope}
      \end{tikzpicture}\\[2pt]
      {\footnotesize \textcolor{goldacc}{$y=x^2,\ x\ge0$} \; and \; \textcolor{forest}{$y=\sqrt{x}$}}
    \end{column}
    \begin{column}{0.54\textwidth}
      \begin{itemize}\setlength{\itemsep}{4pt}
        \item Squaring and square-rooting \textbf{undo} each other --- they are \textbf{inverses}.
        \item Reflecting over $y=x$ \textbf{swaps coordinates}, so domain and range trade places.
        \item \textbf{That is why} the domain of $\sqrt{x}$ is $[0,\infty)$: it is the \emph{range}
              of $x^2$.
      \end{itemize}
      \vspace{2pt}
      \begin{block}{Why only half the parabola?}
        $(-3)^2=3^2=9$ --- two inputs, one output. Reflecting all of it would fail the vertical line
        test, so $y=x^2$ needs a \textbf{restricted domain}. \\[3pt]
        $y=x^3$ never repeats an output, so \emph{all} of it reflects --- and $\sqrt[3]{x}$ needs no
        restriction.
      \end{block}
    \end{column}
  \end{columns}
\end{frame}

% FULL FEATURE READ
\begin{frame}[t, plain]
  \forestheader{Reading the whole feature list}
  \sectionlabel[goldacc]{$f(x)=\sqrt{x+4}-1$}
  \begin{columns}[T]
    \begin{column}{0.42\textwidth}
      \centering
      \begin{tikzpicture}[scale=0.34]
        \lgrid{-6}{7}{-3}{4}
        \begin{scope}
          \clip (-6,-3) rectangle (7,4);
          \draw[very thick, forest, domain=-4:7, samples=100, smooth] plot (\x,{sqrt((\x)+4)-1});
        \end{scope}
        \fill[forest] (-4,-1) circle (6pt);
        \fill[charcoal] (-3,0) circle (4pt);
        \fill[charcoal] (0,1) circle (4pt);
        \fill[charcoal] (5,2) circle (4pt);
      \end{tikzpicture}
    \end{column}
    \begin{column}{0.56\textwidth}
      \footnotesize
      \begin{itemize}\setlength{\itemsep}{3pt}
        \item Domain $[-4,\infty)$ \; (from $x+4\ge0$); \; Range $[-1,\infty)$
        \item \textbf{Endpoint} $(-4,-1)$ --- domain and range both \emph{start} at its coordinates
        \item Zero $(-3,0)$; \; $y$-intercept $(0,1)$
        \item Increasing on $[-4,\infty)$ --- the entire domain
        \item $f(x)<0$ on $[-4,-3)$; \; $f(x)>0$ on $(-3,\infty)$ \\
              {\scriptsize (no interval may leak outside the domain)}
        \item Absolute \textbf{minimum} $-1$ at $x=-4$; no absolute maximum
        \item $x\to+\infty \Rightarrow f(x)\to+\infty$; \; on the left, the graph \textbf{stops}
      \end{itemize}
    \end{column}
  \end{columns}
\end{frame}

% ROADMAP
\begin{frame}[t, plain]
  \forestheader{Where Unit 6 goes next}
  \sectionlabel{Connections}
  \textbf{Today:} read domain (from the radicand), range, endpoint, intercepts, intervals, extrema,
  and end behavior off both radical parents --- and meet the inverse relationship that produced
  them.\\[8pt]
  \begin{itemize}\setlength{\itemsep}{3pt}
    \item \textbf{6.1} $n$th roots \& rational exponents --- what an index really is.
    \item \textbf{6.2--6.3} Simplifying, adding, multiplying, dividing \& rationalizing radicals.
    \item \textbf{6.4} Graphing radical functions \& transformations (the \emph{endpoint} moves).
    \item \textbf{6.5} Solving radical equations \& extraneous solutions.
    \item \textbf{6.6} Composition of functions.
    \item \textbf{6.7} Inverse functions --- today's last idea, made into a method.
  \end{itemize}
\end{frame}

\end{document}
